\chapter{非負行列と Hilbert 射影計量}
\label{ch:found-nonneg-matrix}

\section{Hilbert 射影計量}
\label{sec:sem-hilbert-metric}

Sinkhorn 反復は定数倍の自由度をもつため，収束は射影空間
$\Rpos^n / {\sim}$（$\mathbf{x} \sim \mathbf{x}' \iff \mathbf{x}' = \lambda\mathbf{x},\ \exists\lambda>0$）
の上で測るのが自然である．この空間上の距離を導入する．

\begin{definition}{Hilbert 射影計量}{sem-hilbert-metric}
  正ベクトル $\mathbf{x}, \mathbf{x}' \in \Rpos^n$ に対して，\textbf{Hilbert 射影計量}を
  \[
    \dHil(\mathbf{x}, \mathbf{x}')
    \defeq
    \log \max_{i, k} \frac{x_i\, x'_k}{x_k\, x'_i}
    = \log \frac{\displaystyle\max_i (x_i / x'_i)}{\displaystyle\min_k (x_k / x'_k)}
  \]
  で定める．
\end{definition}

\begin{claim}{射影計量の基本性質}{sem-hilbert-metric-properties}
  $\dHil$ は次を満たす：
  \begin{enumerate}
    \item $\dHil(\mathbf{x}, \mathbf{x}') \geq 0$，かつ等号成立は $\mathbf{x}' = \lambda\mathbf{x}$（ある $\lambda>0$）と同値．
    \item $\dHil(\lambda\mathbf{x}, \mu\mathbf{x}') = \dHil(\mathbf{x}, \mathbf{x}')$（$\lambda, \mu > 0$）．
    \item 対称性 $\dHil(\mathbf{x}, \mathbf{x}') = \dHil(\mathbf{x}', \mathbf{x})$．
    \item 三角不等式 $\dHil(\mathbf{x}, \mathbf{z}) \leq \dHil(\mathbf{x}, \mathbf{y}) + \dHil(\mathbf{y}, \mathbf{z})$．
  \end{enumerate}
  ゆえに $\dHil$ は射影空間 $\Rpos^n / {\sim}$ 上の距離である．
  さらに正規化スライス $\simplex_n^\circ \defeq \{\mathbf{x} \in \Rpos^n : \sum_i x_i = 1\}$ を
  各同値類の代表とすれば，$(\simplex_n^\circ, \dHil)$ は\textbf{完備}距離空間となる．
\end{claim}

\begin{proof}
  $\mathrm{M}(\mathbf{x}/\mathbf{x}') \defeq \max_i (x_i/x'_i)$,
  $\mathrm{m}(\mathbf{x}/\mathbf{x}') \defeq \min_i (x_i/x'_i)$ とおくと
  $\dHil(\mathbf{x}, \mathbf{x}') = \log\bigl(\mathrm{M}(\mathbf{x}/\mathbf{x}') / \mathrm{m}(\mathbf{x}/\mathbf{x}')\bigr)$．

  \textbf{(1)} 常に $\mathrm{M} \geq \mathrm{m}$ だから $\dHil \geq 0$．等号は $\mathrm{M} = \mathrm{m}$，
  すなわち全 $i$ で $x_i/x'_i$ が一定値 $\lambda$ をとること，つまり $\mathbf{x} = \lambda\mathbf{x}'$ と同値．

  \textbf{(2)} $x_i/(\mu x'_i)$ や $(\lambda x_i)/x'_i$ では比が定数倍されるだけで
  $\mathrm{M}/\mathrm{m}$ は不変．

  \textbf{(3)} $\max_{i,k}\frac{x_i x'_k}{x_k x'_i}$ は $(i,k)$ の入れ替えで
  $\max_{i,k}\frac{x'_i x_k}{x'_k x_i}$ に等しい（添字を付け替えるだけ）から対称．

  \textbf{(4)} 任意の $i$ で $\frac{x_i}{z_i} = \frac{x_i}{y_i}\cdot\frac{y_i}{z_i}
  \leq \mathrm{M}(\mathbf{x}/\mathbf{y})\,\mathrm{M}(\mathbf{y}/\mathbf{z})$ だから
  $\mathrm{M}(\mathbf{x}/\mathbf{z}) \leq \mathrm{M}(\mathbf{x}/\mathbf{y})\,\mathrm{M}(\mathbf{y}/\mathbf{z})$．
  同様に $\mathrm{m}(\mathbf{x}/\mathbf{z}) \geq \mathrm{m}(\mathbf{x}/\mathbf{y})\,\mathrm{m}(\mathbf{y}/\mathbf{z})$．
  よって
  \[
    \dHil(\mathbf{x}, \mathbf{z})
    = \log\frac{\mathrm{M}(\mathbf{x}/\mathbf{z})}{\mathrm{m}(\mathbf{x}/\mathbf{z})}
    \leq \log\frac{\mathrm{M}(\mathbf{x}/\mathbf{y})\,\mathrm{M}(\mathbf{y}/\mathbf{z})}
                  {\mathrm{m}(\mathbf{x}/\mathbf{y})\,\mathrm{m}(\mathbf{y}/\mathbf{z})}
    = \dHil(\mathbf{x}, \mathbf{y}) + \dHil(\mathbf{y}, \mathbf{z}).
  \]

  \textbf{完備性．}$(\mathbf{x}^{(\ell)})$ を $\simplex_n^\circ$ 上の $\dHil$-Cauchy 列とする．
  $\dHil$ が小さいことは比 $x_i/x'_i$ が $1$ の近くに揃うことを意味するので，ある $N$ 以降
  $\dHil(\mathbf{x}^{(\ell)}, \mathbf{x}^{(N)}) \leq 1$，すなわち全 $i, \ell\geq N$ で
  $x_i^{(\ell)} / x_i^{(N)} \in [c, C]$（$C/c \leq e$）が成り立つ．
  ゆえに列は $\simplex_n^\circ$ のコンパクト部分集合
  $\{\mathbf{x} : c\,x_i^{(N)} \leq x_i \leq C\,x_i^{(N)},\ \sum_i x_i = 1\}$ に入り，
  収束部分列をもち，その極限 $\mathbf{x}^\star$ は全成分正である．
  $\dHil$ は $\Rpos^n$ 上で連続だから，Cauchy 性と合わせて
  $\dHil(\mathbf{x}^{(\ell)}, \mathbf{x}^\star) \to 0$．よって完備である．
\end{proof}

%% ============================================================
%% 4.3 Birkhoff の縮小定理
%% ============================================================
\section{Birkhoff の縮小定理}
\label{sec:sem-birkhoff}

本節の目標は，正行列 $\mathbf{K}$ が線形写像 $\mathbf{x} \mapsto \mathbf{K}\mathbf{x}$ として
Hilbert 射影計量を\textbf{狭義に縮小}し，その縮小率が成分の歪みだけで定まる
明示的な定数 $\lambda(\mathbf{K}) < 1$ で与えられることを示すことである．

まず，正行列による像の成分比が，入力の成分比の\textbf{重み付き平均}になるという基本観察から始める．

\begin{claim}{像の成分比は重み付き平均}{sem-weighted-average}
  $\mathbf{x}, \mathbf{y} \in \Rpos^m$ とし $t_k \defeq x_k / y_k$ とおく．
  $\mathbf{K} \in \Rpos^{n\times m}$ に対し，重み
  $w_{i,k} \defeq \dfrac{K_{i,k} y_k}{(\mathbf{K}\mathbf{y})_i}$ は
  $w_{i,k} > 0$, $\sum_k w_{i,k} = 1$ を満たし，
  \[
    \frac{(\mathbf{K}\mathbf{x})_i}{(\mathbf{K}\mathbf{y})_i}
    = \sum_{k} w_{i,k}\, t_k
  \]
  が成り立つ．
\end{claim}

\begin{proof}
  $w_{i,k} > 0$ と $\sum_k w_{i,k} = \sum_k K_{i,k}y_k / (\mathbf{K}\mathbf{y})_i = 1$ は定義から明らか．
  $x_k = t_k y_k$ より
  $(\mathbf{K}\mathbf{x})_i = \sum_k K_{i,k} t_k y_k$ だから，
  $(\mathbf{K}\mathbf{x})_i / (\mathbf{K}\mathbf{y})_i = \sum_k \frac{K_{i,k}y_k}{(\mathbf{K}\mathbf{y})_i} t_k
  = \sum_k w_{i,k} t_k$．
\end{proof}

次に，像 $\mathbf{K}\mathbf{x}$ の取りうる範囲の「射影的な広がり」（射影直径）が
成分比の最大歪みで定まることを示す．

\begin{claim}{射影直径}{sem-projective-diameter}
  $\mathbf{K} \in \Rpos^{n\times m}$ の\textbf{射影直径}を
  $\Delta(\mathbf{K}) \defeq \sup_{\mathbf{x}, \mathbf{y} \in \Rpos^m} \dHil(\mathbf{K}\mathbf{x}, \mathbf{K}\mathbf{y})$
  とおくと，
  \[
    \Delta(\mathbf{K}) = \log \eta(\mathbf{K}),
    \qquad
    \eta(\mathbf{K}) \defeq \max_{i,j,k,\ell} \frac{K_{i,k}\, K_{j,\ell}}{K_{j,k}\, K_{i,\ell}}.
  \]
\end{claim}

\begin{proof}
  \textbf{上界．}正数 $a_k, b_k > 0$ に対する中間値不等式
  $\frac{\sum_k a_k}{\sum_k b_k} \leq \max_k \frac{a_k}{b_k}$（分母を払えば明らか）を用いる．
  任意の $\mathbf{x} > 0$ と行 $i, j$ に対し
  \[
    \frac{(\mathbf{K}\mathbf{x})_i}{(\mathbf{K}\mathbf{x})_j}
    = \frac{\sum_k K_{i,k} x_k}{\sum_k K_{j,k} x_k}
    \leq \max_k \frac{K_{i,k}}{K_{j,k}}.
  \]
  ゆえに $\mathbf{x}, \mathbf{y} > 0$ に対し
  \[
    \frac{(\mathbf{K}\mathbf{x})_i\,(\mathbf{K}\mathbf{y})_j}{(\mathbf{K}\mathbf{x})_j\,(\mathbf{K}\mathbf{y})_i}
    \leq \max_k \frac{K_{i,k}}{K_{j,k}} \cdot \max_\ell \frac{K_{j,\ell}}{K_{i,\ell}}
    = \max_{k,\ell} \frac{K_{i,k}\,K_{j,\ell}}{K_{j,k}\,K_{i,\ell}}
    \leq \eta(\mathbf{K}).
  \]
  $i, j$ の最大をとって $\dHil(\mathbf{K}\mathbf{x}, \mathbf{K}\mathbf{y}) \leq \log\eta(\mathbf{K})$，
  よって $\Delta(\mathbf{K}) \leq \log\eta(\mathbf{K})$．

  \textbf{下界．}$\eta(\mathbf{K})$ を達成する添字を $(i^*, j^*, k^*, \ell^*)$ とする．
  $\mathbf{x} = \mathbf{e}_{k^*} + \delta\ones_m$, $\mathbf{y} = \mathbf{e}_{\ell^*} + \delta\ones_m$（$\delta > 0$）とおくと
  $\mathbf{x}, \mathbf{y} > 0$ で，$\delta \to 0+$ のとき
  $\mathbf{K}\mathbf{x} \to \mathbf{K}_{\cdot,k^*}$（第 $k^*$ 列），$\mathbf{K}\mathbf{y} \to \mathbf{K}_{\cdot,\ell^*}$ となり
  \[
    \dHil(\mathbf{K}\mathbf{x}, \mathbf{K}\mathbf{y})
    \to \log\max_{i,j}\frac{K_{i,k^*}K_{j,\ell^*}}{K_{j,k^*}K_{i,\ell^*}}
    = \log\eta(\mathbf{K}).
  \]
  ゆえに $\Delta(\mathbf{K}) \geq \log\eta(\mathbf{K})$．以上より等号が成り立つ．
\end{proof}

\begin{theorem}{Birkhoff の縮小定理}{sem-birkhoff}
  $\mathbf{K} \in \Rpos^{n\times m}$ は Hilbert 射影計量を縮小する：すべての
  $\mathbf{x}, \mathbf{y} \in \Rpos^m$ に対し
  \[
    \dHil(\mathbf{K}\mathbf{x}, \mathbf{K}\mathbf{y})
    \leq \lambda(\mathbf{K})\, \dHil(\mathbf{x}, \mathbf{y}),
    \qquad
    \lambda(\mathbf{K})
    \defeq \tanh\!\left(\frac{\Delta(\mathbf{K})}{4}\right)
    = \frac{\sqrt{\eta(\mathbf{K})} - 1}{\sqrt{\eta(\mathbf{K})} + 1} \in [0, 1).
  \]
  $\eta(\mathbf{K}) < \infty$ ゆえ $\lambda(\mathbf{K}) < 1$（狭義縮小）である．
\end{theorem}

\begin{proof}
  $\eta \defeq \eta(\mathbf{K})$, $s \defeq e^{\dHil(\mathbf{x}, \mathbf{y})} \geq 1$ とおく．
  $s$ は成分比 $t_k = x_k/y_k$ の最大歪み $M/m$ に等しい（$M = \max_k t_k$, $m = \min_k t_k$）．
  $s = 1$ なら $\mathbf{x} \sim \mathbf{y}$ で両辺 $0$ ゆえ自明．以下 $s > 1$ とする．

  \textbf{Step 1（二行・二値への縮約）．}
  主張~\ref{clm:sem-weighted-average} より $\phi_i \defeq (\mathbf{K}\mathbf{x})_i/(\mathbf{K}\mathbf{y})_i
  = \sum_k w_{i,k} t_k$ であり，
  $\dHil(\mathbf{K}\mathbf{x}, \mathbf{K}\mathbf{y}) = \log\bigl(\max_i \phi_i / \min_i \phi_i\bigr)$．
  最大・最小を達成する行を $i^*, j^*$ とすれば，評価すべきは $\phi_{i^*}/\phi_{j^*}$ である．
  $\phi_{i^*}/\phi_{j^*} = \bigl(\sum_k w_{i^*,k} t_k\bigr)/\bigl(\sum_k w_{j^*,k} t_k\bigr)$ を
  各 $t_k$ の関数と見ると，$t_k$ について一次分数関数（Möbius 型）であり，
  分母は正だから $t_k$ に関して単調で，端点 $t_k \in \{m, M\}$ で最大化される．
  よって最悪値の評価では各 $t_k$ は $m$ か $M$ のいずれかとしてよい．

  列を $H \defeq \{k : t_k = M\}$, $L \defeq \{k : t_k = m\}$ に分け，
  $p \defeq \sum_{k\in H} w_{i^*,k}$, $q \defeq \sum_{k\in H} w_{j^*,k}$ とおく（$p, q \in [0,1]$）と
  \[
    \frac{\phi_{i^*}}{\phi_{j^*}} = \frac{pM + (1-p)m}{qM + (1-q)m}.
  \]
  これを最大化するので $p$ を大きく $q$ を小さくとる向きに考える．

  \textbf{Step 2（重みの歪み制約）．}
  重みの定義より任意の行 $i, j$ と列 $k, \ell$ で
  \[
    \frac{w_{i,k}\, w_{j,\ell}}{w_{i,\ell}\, w_{j,k}}
    = \frac{K_{i,k} y_k \cdot K_{j,\ell} y_\ell}{K_{i,\ell} y_\ell \cdot K_{j,k} y_k}
    = \frac{K_{i,k}\, K_{j,\ell}}{K_{i,\ell}\, K_{j,k}}
    \leq \eta
  \]
  （$y$ は約分される）．$k \in H$, $\ell \in L$ にわたって和をとると
  \[
    \Bigl(\sum_{k\in H} w_{i^*,k}\Bigr)\Bigl(\sum_{\ell\in L} w_{j^*,\ell}\Bigr)
    \leq \eta \Bigl(\sum_{\ell\in L} w_{i^*,\ell}\Bigr)\Bigl(\sum_{k\in H} w_{j^*,k}\Bigr),
    \quad\text{すなわち}\quad
    p(1-q) \leq \eta\,(1-p)\,q.
  \]
  オッズ $\displaystyle u \defeq \frac{p}{1-p}$, $\displaystyle v \defeq \frac{q}{1-q}$ で書けば $u \leq \eta v$．

  \textbf{Step 3（一変数最適化）．}
  $p M + (1-p)m = (1-p)(uM + m)$, $\frac{1-p}{1-q} = \frac{1+v}{1+u}$ を使うと，
  比の最大化は $u = \eta v$（制約の境界）で達成され，$v > 0$ を動かして
  \[
    \frac{\phi_{i^*}}{\phi_{j^*}}
    \leq f(v) \defeq \frac{(1+v)(\eta s v + 1)}{(1+\eta v)(s v + 1)}
  \]
  を最大化すればよい（両辺を $m$ で割り $s = M/m$ を用いた）．
  $f(v) \to 1$（$v \to 0+$ および $v \to \infty$）である．$f$ は分母が $(0,\infty)$ 上で正ゆえ
  $(0,\infty)$ 上で微分可能（定義~\ref{def:sem-differentiable}）であり，停留点を求めるため対数微分すると
  \[
    \frac{f'(v)}{f(v)}
    = \frac{1}{1+v} + \frac{\eta s}{\eta s v + 1}
      - \frac{\eta}{1+\eta v} - \frac{s}{s v + 1}.
  \]
  右辺の第 1・第 4 項，第 2・第 3 項をそれぞれ通分すると
  \[
    \frac{1}{1+v} - \frac{s}{s v + 1}
    = \frac{1-s}{(1+v)(s v + 1)},
    \qquad
    \frac{\eta s}{\eta s v + 1} - \frac{\eta}{1+\eta v}
    = \frac{\eta(s-1)}{(\eta s v + 1)(1+\eta v)}
  \]
  だから
  \[
    \frac{f'(v)}{f(v)}
    = (s-1)\left[
        \frac{\eta}{(\eta s v + 1)(1+\eta v)}
        - \frac{1}{(1+v)(s v + 1)}
      \right].
  \]
  $s > 1$ ゆえ $f'(v) = 0$ は角括弧内が $0$ となること，すなわち
  $\eta\,(1+v)(s v + 1) = (\eta s v + 1)(1+\eta v)$ と同値である．両辺を展開すると
  \[
    \eta s\,v^2 + \eta(s+1)\,v + \eta
    = \eta^2 s\,v^2 + \eta(s+1)\,v + 1
  \]
  となり，一次の項 $\eta(s+1)v$ が相殺して $\eta s\,v^2 (1-\eta) = 1 - \eta$，
  $\eta > 1$ より $\eta s\,v^2 = 1$．したがって $(0,\infty)$ における唯一の停留点は
  $v^* = 1/\sqrt{\eta s}$ である．$f(v) \to 1$（両端点）かつ
  $f(v^*) > 1$（後述）だから，$v^*$ が最大値を与える．これを代入すると
  \[
    1 + v^* = \frac{\sqrt{\eta s}+1}{\sqrt{\eta s}},\quad
    \eta s\,v^* + 1 = \sqrt{\eta s}+1,\quad
    1 + \eta v^* = \frac{\sqrt{\eta}+\sqrt{s}}{\sqrt{s}},\quad
    s\,v^* + 1 = \frac{\sqrt{\eta}+\sqrt{s}}{\sqrt{\eta}}
  \]
  （$\sqrt{\eta}\sqrt{s} = \sqrt{\eta s}$ を用いた）であり，
  \[
    \max_{v>0} f(v) = f(v^*)
    = \frac{(1+v^*)(\eta s\,v^* + 1)}{(1+\eta v^*)(s\,v^* + 1)}
    = \left(\frac{\sqrt{\eta s} + 1}{\sqrt{\eta} + \sqrt{s}}\right)^{\!2}.
  \]
  ここで $\sqrt{\eta s} + 1 > \sqrt{\eta} + \sqrt{s}$
  （$\iff (\sqrt{\eta}-1)(\sqrt{s}-1) > 0$，$\eta, s > 1$ より成立）から
  $f(v^*) > 1$ である．
  ゆえに
  \[
    \dHil(\mathbf{K}\mathbf{x}, \mathbf{K}\mathbf{y})
    = \log\frac{\phi_{i^*}}{\phi_{j^*}}
    \leq 2\log\frac{\sqrt{\eta s} + 1}{\sqrt{\eta} + \sqrt{s}}.
  \]

  \textbf{Step 4（率の評価）．}
  $\theta \defeq \tfrac12\log s \geq 0$, $\beta \defeq \tfrac12\log\eta \geq 0$ とおくと
  $\sqrt{s} = e^\theta$, $\sqrt{\eta} = e^\beta$ で，
  $\dHil(\mathbf{x}, \mathbf{y}) = \log s = 2\theta$,
  $\lambda(\mathbf{K}) = \tanh(\tfrac14\log\eta) = \tanh(\beta/2)$ である．
  示すべき不等式は
  \[
    \psi(\theta) \defeq
    \log\bigl(e^{\beta+\theta} + 1\bigr) - \log\bigl(e^\beta + e^\theta\bigr)
    - \tanh(\beta/2)\,\theta \leq 0
    \qquad (\theta \geq 0)
  \]
  に帰着する（左辺の最初の二項が
  $\log\frac{\sqrt{\eta s}+1}{\sqrt\eta+\sqrt s}$，全体を $2$ 倍すると上の評価）．
  $\psi(0) = 0$ である．また
  \[
    \psi'(\theta)
    = \frac{e^{\beta+\theta}}{e^{\beta+\theta}+1}
    - \frac{e^\theta}{e^\beta + e^\theta}
    - \tanh(\beta/2)
  \]
  で，$\psi'(0) = \frac{e^\beta - 1}{e^\beta + 1} - \tanh(\beta/2) = 0$
  （$\frac{e^\beta-1}{e^\beta+1} = \tanh(\beta/2)$）．さらに
  \[
    \psi''(\theta)
    = \frac{e^{\beta+\theta}}{(e^{\beta+\theta}+1)^2}
    - \frac{e^{\beta+\theta}}{(e^\beta + e^\theta)^2}
    = e^{\beta+\theta}\left[\frac{1}{(e^{\beta+\theta}+1)^2} - \frac{1}{(e^\beta+e^\theta)^2}\right].
  \]
  ここで
  $(e^{\beta+\theta} + 1) - (e^\beta + e^\theta) = (e^\beta - 1)(e^\theta - 1) \geq 0$
  ($\beta, \theta \geq 0$) より $(e^{\beta+\theta}+1)^2 \geq (e^\beta + e^\theta)^2$，ゆえに $\psi'' \leq 0$．
  したがって $\psi$ は $\theta \geq 0$ で凹で，$\psi(0) = \psi'(0) = 0$ だから $\psi(\theta) \leq 0$．
  これは
  \[
    \dHil(\mathbf{K}\mathbf{x}, \mathbf{K}\mathbf{y})
    \leq 2\log\frac{\sqrt{\eta s}+1}{\sqrt\eta+\sqrt s}
    \leq \tanh(\beta/2)\cdot 2\theta
    = \lambda(\mathbf{K})\,\dHil(\mathbf{x}, \mathbf{y})
  \]
  を意味する．最後に $\eta = e^{2\beta} < \infty$ より $\lambda(\mathbf{K}) = \tanh(\beta/2) < 1$．
\end{proof}

\begin{remark}{率は $s\to1$ で達成される}{sem-birkhoff-tightness}
  Step 4 で $\psi$ が凹で $\psi(0)=\psi'(0)=0$ だったことは，比
  $\dHil(\mathbf{K}\mathbf{x},\mathbf{K}\mathbf{y})/\dHil(\mathbf{x},\mathbf{y})$ の上限 $\lambda(\mathbf{K})$ が
  $s \to 1+$（入力がほぼ比例）で漸近的に達成されることを示す．一方 $s \to \infty$ では
  $\dHil(\mathbf{K}\mathbf{x},\mathbf{K}\mathbf{y}) \to \Delta(\mathbf{K}) = \log\eta$ で頭打ちになる
  （像の射影直径は有界）．Gibbs カーネル $K_{i,j} = e^{-C_{i,j}/\varepsilon}$ では
  $\eta(\mathbf{K}) = \exp\bigl(\tfrac{1}{\varepsilon}\max_{i,j,k,\ell}(C_{j,k}+C_{i,\ell}-C_{i,k}-C_{j,\ell})\bigr)
  \leq e^{2\norm{\mathbf{C}}_\infty/\varepsilon}$ となり，$\varepsilon$ が小さいほど
  $\eta \to \infty$, $\lambda(\mathbf{K}) \to 1$ で縮小は弱くなる．
\end{remark}

%% ============================================================
%% 4.4 Sinkhorn の線形収束
%% ============================================================
